% ============================================================================
\section{Oversight, monitoring, and evaluation}
% ============================================================================

\begin{frame}{Logging and Monitoring Agent Actions}
  \begin{itemize}\itemsep0.5em
    \item[\textcolor{KAred}{$\blacktriangleright$}] Because an agent acts, its actions can be audited in a way a chatbot's reasoning cannot. Log every tool call, argument, and result.
    \item[\textcolor{KAred}{$\blacktriangleright$}] \textbf{Anomaly detection:} flag behavior that departs from the task, such as an agent asked to summarize email that suddenly sends one, or reads files it was never pointed at.
    \item[\textcolor{KAred}{$\blacktriangleright$}] Monitoring depends on independence. If the same model that was hijacked also writes the summary of what it did, the log inherits the compromise.
    \item[\textcolor{KAred}{$\blacktriangleright$}] Prefer monitors that watch the actions themselves, such as tool calls and network egress, over the agent's self-report.
  \end{itemize}
\end{frame}

\begin{frame}{Red-Teaming and Automated Attack Discovery}
  \begin{itemize}\itemsep0.5em
    \item[\textcolor{KAred}{$\blacktriangleright$}] \textbf{Red-teaming:} deliberately attack your own agent by planting injections, crafting malicious pages, and probing tool misuse before an adversary does.
    \item[\textcolor{KAred}{$\blacktriangleright$}] Manual red-teaming does not scale to every page an agent might read, so the attack search is increasingly automated: generate and mutate injections, then keep what works.
    \item[\textcolor{KAred}{$\blacktriangleright$}] The persuasion-and-manipulation grid that TRAP uses to attack also serves as a test suite: run it against a candidate agent and measure the hijack rate.
    \item[\textcolor{KAred}{$\blacktriangleright$}] Robustness is best treated as a number tracked across releases, not a one-time check.
  \end{itemize}
\end{frame}

\begin{frame}{Benchmarks for Agent Harm and Robustness}
  \begin{itemize}\itemsep0.5em
    \item[\textcolor{KAred}{$\blacktriangleright$}] \textbf{AgentHarm} measures misuse: will the agent execute an explicitly malicious multi-step task? It has 110 tasks across 11 harm categories.
    \item[\textcolor{KAred}{$\blacktriangleright$}] \textbf{AgentDojo} measures adversarial robustness with a dynamic environment that pairs benign user tasks with injection attacks and candidate defenses.
    \item[\textcolor{KAred}{$\blacktriangleright$}] \textbf{TRAP} measures indirect injection on web agents inside high-fidelity clones of six popular sites, standing in for Amazon, Gmail, Calendar, LinkedIn, DoorDash, and Upwork.
    \item[\textcolor{KAred}{$\blacktriangleright$}] Benchmarks cover the attacks we already know. A passing score is a lower bound on safety, not a guarantee.
  \end{itemize}
\end{frame}

\begin{frame}{The Limits of Agent Benchmarks}
  \begin{itemize}\itemsep0.5em
    \item[\textcolor{KAred}{$\blacktriangleright$}] Most web-agent benchmarks feed the model the page's underlying structure as text: its \textbf{DOM} (Document Object Model) or accessibility tree. This is convenient but skips the vision-heavy, specialized interfaces real users face.
    \item[\textcolor{KAred}{$\blacktriangleright$}] A benchmark can therefore look near-saturated while the agent is still blind to anything that is not clean structured text.
    \item[\textcolor{KAred}{$\blacktriangleright$}] On GauntletBench, a suite of vision-intensive tasks in five specialized professional applications, non-expert humans succeed on over 80\% of tasks while the strongest agent reaches 19\%.
    \item[\textcolor{KAred}{$\blacktriangleright$}] For safety this cuts both ways: today's agents are less capable than headlines suggest, but evaluation that overstates capability also \textbf{understates the failure modes} we must guard against.
  \end{itemize}
  \imgsrc{Vysotskyi et al., ``Running the Gauntlet: Re-evaluating the Capabilities of Agents Beyond Familiar Environments'' (GauntletBench), \href{https://arxiv.org/abs/2606.14397}{arXiv:2606.14397}.}
\end{frame}

\begin{frame}{Incident Response for Deployed Agents}
  \begin{itemize}\itemsep0.5em
    \item[\textcolor{KAred}{$\blacktriangleright$}] Assume something will go wrong in production, and plan the response in advance.
    \item[\textcolor{KAred}{$\blacktriangleright$}] \textbf{Contain:} a kill switch that halts the agent and revokes its credentials immediately.
    \item[\textcolor{KAred}{$\blacktriangleright$}] \textbf{Investigate:} action logs detailed enough to reconstruct what happened and what was affected.
    \item[\textcolor{KAred}{$\blacktriangleright$}] \textbf{Recover and learn:} reverse what can be reversed, notify affected parties, and turn the root cause into a new red-team test so it cannot recur silently.
  \end{itemize}
\end{frame}
